\documentclass[10pt]{mypackage}

\usepackage[uprightgreeks,light]{kpfonts}
\renewcommand{\coloneq}{\coloneqq}
\renewcommand{\D}{\mathbb{D}}
%\usepackage{homework}
\usepackage{notes}

\usepackage[ backend=bibtex, style = alphabetic, sorting=ynt ]{biblatex}
\addbibresource{all_references.bib}

\usepackage{parskip}

\fancyhf{}
\fancyhead[R]{Avinash Iyer}
\fancyhead[L]{Measure-Preserving Actions and Ergodicity}
\fancyfoot[C]{\thepage}

\setcounter{secnumdepth}{0}

\begin{document}
\RaggedRight
\section{Preliminary: Positive Definite Functions, the GNS Construction, and Representations}%
We discuss some preliminaries on representations of groups via the GNS construction.
\begin{definition}
  A function $f\colon G\rightarrow \C$ is said to be \textit{positive definite} if, for all $s_1,\dots,s_n\in G$ and $c_1,\dots,c_n\in \C$, one has
  \begin{align*}
    \sum_{i,j=1}^{n} \overline{c_i}c_j f\left( s_i^{-1}s_j \right) &\geq 0.
  \end{align*}
  That is, the matrix $A = \left( f\left( s_i^{-1}s_j \right) \right)_{i,j}$ is a positive operator on $\C^n$. We let $\mathcal{P}(G)$ denote the set of all positive definite functions on $G$.
\end{definition}
\begin{proposition}
  Let $f\colon G\rightarrow \C$ be a positive definite function. Then,
  \begin{enumerate}[(i)]
    \item $f\left( s^{-1} \right) = \overline{f(s)}$ for all $s\in G$;
    \item $ \left\vert f\left( s \right) \right\vert\leq f\left( e \right) $ for all $s\in G$.
  \end{enumerate}
\end{proposition}
\begin{proof}
  Let $s\in G$. If we let $s_1 = e$ and $s_2 = s$, then the matrix
  \begin{align*}
    A &= \begin{pmatrix}f(e) & f(s) \\ f\left( s^{-1} \right) & f\left( e \right)\end{pmatrix}
  \end{align*}
  is positive. Therefore, $A$ is Hermitian, so $ f\left( s^{-1} \right) = \overline{f\left( s \right)} $.

  Moreover, if
  \begin{align*}
    \xi &= \begin{pmatrix}1\\ -\frac{ \overline{f\left( s \right)} }{\left\vert f\left( s \right) \right\vert}\end{pmatrix},
  \end{align*}
  then we deduce from the fact that $ \iprod{A\xi}{\xi} \geq 0$ that $\left\vert f\left( s \right) \right\vert\leq f\left( e \right)$.
\end{proof}
We let $\mathcal{P}_1(G)$ denote the set of all positive definite functions $\varphi\colon G\rightarrow \C$ with $\varphi\left( e \right) = 1$. The weak* topology on $\mathcal{P}_1\left( G \right)$ emerges by regarding it as a subset of the closed unit ball of $\ell^{\infty}\left( G \right)$, which is the dual of $\ell^1\left( G \right)$.

On the closed unit ball of $\ell^{\infty}\left( G \right)$, the weak* topology and the topology of pointwise convergence coincide.

If $\pi\colon G\rightarrow B\left( H \right)$ is a unitary representation, and $\xi\in H$, then we may define the function
\begin{align*}
  f(s) &= \iprod{\pi(s)\xi}{\xi}.
\end{align*}
Observe that
\begin{align*}
  \sum_{i,j=1}^{n} \overline{c_i}c_j f\left( s_i^{-1}s_j \right) &= \iprod{\sum_{i=1}^{n}c_i\pi\left( s_i \right)\xi}{\sum_{i=1}^{n}c_i\pi\left( s_i \right)\xi}\\
                                                                 &\geq 0,
\end{align*}
so that $f$ is positive definite.

One can see a parallel between the case of positive definite functions and the case of states on $C^{\ast}$-algebras. This parallel extends further.
\begin{theorem}[GNS Construction]
  Let $f\colon G\rightarrow \C$ be a positive definite function. There is a unitary representation $\pi\colon G\rightarrow B\left( H \right)$ and a cyclic vector $\eta\in H$ such that
  \begin{align*}
    f(s) &= \iprod{\pi(s)\eta}{\eta}.
  \end{align*}
  In this case, $\norm{f}_{\ell^{\infty}} = \norm{\eta}^2$.

  The triple $\left( \pi,H,\eta \right)$ is unique in the sense that if $\pi'\colon G\rightarrow B\left( H' \right)$ is another unitary representation, and $\eta'$ is a cyclic vector in $H'$ such that $f(s) = \iprod{\pi'(s)\eta'}{\eta'}$ for all $s\in G$, then there is a unitary isomorphism from $H$ to $H'$ that intertwines $\pi$ and $\pi'$ and maps $\eta\mapsto \eta'$.
\end{theorem}
\begin{proof}
  On $\C[G]$, the vector space of complex-valued functions on $G$, define the sesquilinear form
  \begin{align*}
    \iprod{\xi}{\zeta} &= \sum_{s,t\in G} \overline{\zeta(t)} \xi(s) f\left( t^{-1}s \right).
  \end{align*}
  This satisfies $ \iprod{\xi}{\xi} \geq 0 $ for all $\xi\in V$ since $f$ is positive-definite.

  Let $ \widetilde{V} $ be the quotient of $V$ by the subspace of vectors $\xi$ satisfying $ \iprod{\xi}{\xi} = 0 $. Then, $ \iprod{\xi}{\xi} $ descends to an inner product on $ \widetilde{V} $, and completing with respect to the induced norm yields the Hilbert space $H$.

  Defining $\left( \pi(s)\xi \right)(t) = \xi\left( s^{-1}t \right)$ for $\xi\in V$ and $s,t\in G$, we get a homomorphism $\pi$ from $G$ into the space of linear transformations on $V$ that satisfies
  \begin{align*}
    \iprod{\pi(s)\xi}{\pi(s)\zeta} &= \sum_{t,u\in G} \overline{\zeta\left( s^{-1}u \right)}\xi\left( s^{-1}t \right)f\left( \left( s^{-1}u \right)^{-1}s^{-1}t \right)\\
                                  &= \sum_{t,u\in G} \overline{\zeta(u)} \xi\left( t \right) f\left( u^{-1}t \right)\\
                                  &= \iprod{\xi}{\zeta}.
  \end{align*}
  Therefore, we have a unitary representation of $G$ on $H$ which we denote by $\pi$. Taking $\eta$ in the theorem statement to be $\delta_e$ in $H$.

  Suppose now that $\pi'\colon G\rightarrow B\left( H' \right)$ is another unitary representation and $\eta'$ another cyclic vector in $H'$ such that $f(s) = \iprod{\pi'(s)\eta'}{\eta'}$ for all $s\in G$. Define a linear map from the linear span of $\pi(G)\eta$ to $H'$ by taking
  \begin{align*}
    \sum_{s\in G} c_s\pi(s)\eta &\mapsto \sum_{s\in G} c_s\pi'(s)\eta',
  \end{align*}
  where the $c_s$ are all but finitely many zero.

  This map is isometric, and extends to a unitary isomorphism $H\rightarrow H'$ which satisfies the desired properties.
\end{proof}
Next, we discuss weak containment of representations.
\begin{definition}
  Let $\pi\colon G\rightarrow B\left( H \right)$ and $\rho\colon G\rightarrow B\left( K \right)$ be unitary representations. We say that $\pi$ is \textit{weakly contained} in $\rho$, and write $ \pi\prec \rho $, if for every finite $F\subseteq G$, $\xi\in H$, and $\ve > 0$, there are $\zeta_1,\dots,\zeta_n\in K$ such that
  \begin{align*}
    \left\vert \iprod{\pi(s)\xi}{\xi} - \sum_{i=1}^{n} \iprod{\rho(s)\zeta_i}{\zeta_i} \right\vert &< \ve.
  \end{align*}
  That is, the matrix coefficients of $\pi(s)$ can be approximated by the matrix coefficients of $\rho(s)$.
\end{definition}
Equivalently, we can say that a positive definite function on $G$ is \textit{associated} to $\pi$ if it is of the form $s\mapsto \iprod{\pi(s)\xi}{\xi}$, and that $\pi\prec \rho$ if any positive definite function associated with $\pi$ is contained in the closure of the set of all finite sums of positive definite functions associated with $\rho$.
\begin{lemma}
  Let $\pi\colon G\rightarrow B\left( H \right)$ and $\rho\colon G\rightarrow B\left( K \right)$ be unitary representations, and let $T\colon H\rightarrow K$ an intertwining operator for $\pi$ and $\rho$. Then, $ \ker\left( T \right)^{\perp} $ and $ \overline{\img\left( T \right)} $ are closed $G$-invariant subspaces of $H$ and $K$ respectively, and the corresponding subrepresentations of $\pi$ and $\rho$ are equivalent.
\end{lemma}
\section{Preliminary: Conditional Expectations}%

\section{Probability Measure-Preserving Actions}%
\begin{definition}
  Let $G\curvearrowright \left( X,\mu \right)$ be a probability measure preserving (pmp) action of $G$. Then, the \textit{Koopman representation} is the unitary representation $\kappa\colon G\rightarrow B\left( L^2\left( X \right) \right)$ given by
  \begin{align*}
    \kappa\left( s \right)f(x) &= f\left( s^{-1}x \right).
  \end{align*}
\end{definition}
Since $\C 1$ is invariant under this action, we will usually restrict our focus to its orthogonal complemented, denoted $L^2\left( X \right)\ominus \C 1$.
\begin{definition}
  A pmp action $G\curvearrowright \left( X,\mu \right)$ is said to be \textit{ergodic} if, for any $G$-invariant subset $A\subseteq X$, $A$ is either null or co-null.
\end{definition}
\begin{definition}
  The action $G\curvearrowright \left( X,\mu \right)$ is said to be (essentially) free if there is a $G$-invariant subset $X_0\subseteq X$ with $\mu\left( X_0 \right) = 1$ and, whenever $sx = x$ for any $x\in X_0$, we have $s = e$.
\end{definition}
We will usually specify that $G$ is a countable discrete group, and $X$ is \href{https://ai-bearing.github.io/Classes_and_Homework/After\%20College/Other\%20Notes/orbit_equivalence.pdf}{standard}.

We can can characterize freeness through the action on subsets, which will allow us to better understand the interactions between ergodicity and freeness.
\begin{proposition}
  The action $G\curvearrowright \left( X,\mu \right)$ is free if and only if for every finite subset $F\subseteq G$ and any non-null set $A\subseteq X$, there is a non-null set $B\subseteq A$ such that $sB\cap tB = \emptyset$ for all distinct $s,t\in F$.
\end{proposition}
\begin{proof}
  If the action is not free, then since $G$ is countable, one of the sets $\set{x\in X | sx = x}$ for some $s\in G\setminus \set{e}$ is positive measure. This set, along with the pair $\set{e,s}$, fails to have the property in the proposition statement.

  In the converse direction, suppose the action is free. We induct on the number of elements in $F$. If $A$ is a non-null subset of $X$, then if $F$ is empty or a singleton, we find the conclusion by taking $B = A$. Suppose $F$ is an arbitrary nonempty finite subset of $G$, and there is a non-null subset $B_0\subseteq A$, with $sB_0\cap S'B_0 = \emptyset$ for all distinct $s,s'\in F$. Let $t\in G\setminus F$, and let $F = \set{s_1,\dots,s_n}$. A result from descriptive set theory provides for $A\supseteq B_1\supseteq\cdots\supseteq B_n$ with $tB_k\cap S_kB_k = \emptyset$ for all $k=1,\dots,n$. The set $B_n$ has the desired property with respect to $F\cup \set{t}$ and $A$.
\end{proof}
\begin{proposition}
  If $G\curvearrowright X$ is a pmp action, then the action is ergodic if and only if for all subsets $A,B\subseteq X$ with positive measure, there is $s\in G$ such that $\mu\left( sA\cap B \right) > 0$.
\end{proposition}
\begin{proof}
  If $A$ and $B$ are subsets of $X$ with positive measure, then the set
  \begin{align*}
    A' &\coloneq \bigcup_{s\in G}sA
  \end{align*}
  satisfies $sA'=A'$ for all $s\in G$, meaning that it must have measure one. In particular, its intersection with $B$ has the same measure as $B$, so since $G$ is countable, at least one of the sets $sA\cap B$ has positive measure.

  In the reverse direction, if $A\subseteq X$ is such that $0 < \mu(A) < 1$, then there is some $s\in G$ such that $\mu\left( sA\cap \left( X\setminus A \right) \right) > 0$, so $A$ fails to be $G$-invariant.
\end{proof}
Working with the definition of ergodicity directly is quite difficult, so we will often use the Koopman representation, and can apply the terminology to general unitary representations.
\begin{definition}
  A unitary representation $\pi\colon G\rightarrow B\left( H \right)$ is called \textit{ergodic} if $\pi(s)\xi = \xi$ for all $s\in G$ implies $\xi = 0$. That is, there are no nonzero invariant vectors.
\end{definition}
\begin{proposition}
  The action $G\curvearrowright \left( X,\mu \right)$ is ergodic if and only if the restriction of the Koopman representation to the orthogonal complement $L^2\left( X,\mu \right)\ominus \C 1$ of the constants is ergodic.
\end{proposition}
\begin{proof}
  If $A$ is $G$-invariant with $0 < \mu(A) < 1$, then $\chi_A-\mu(A)1$ is a nonzero $G$-invariant vector in $L^2\left( X,\mu \right)\ominus \C 1$.

  Conversely, if $f$ is a nonzero $G$-invariant vector in $L^2\left( X,\mu \right)\ominus \C1$, then we may find a measurable subset $D\subseteq \C$ with $0 < \mu\left( f^{-1}\left( D \right) \right) < 1$. For each $s\in G$, the sets $sf^{-1}\left( D \right)\triangle f^{-1}\left( D \right)$ have measure zero, as they are contained in the union over $n\in \N$ of the sets $\set{x\in X | \left\vert f(x)-sf(x) \right\vert > 1/n}$, each of which has measure zero since $\norm{f - sf}_{L^2} = 0$. Thus, the action is not ergodic.
\end{proof}
One of the distinctive features of actions of infinite groups on finite sets is that these actions cannot be free. We can translate this to pmp actions $G\curvearrowright \left( X,\mu \right)$ by replacing points by sets of nonzero measure.

Given $A\subseteq X$ of nonzero measure, there exists some nontrivial $s\in G$ with $\mu\left( sA\cap A \right) > 0$, since otherwise we would have $ sA $ are almost everywhere disjoint for every $s\in G$, whence if $F\subseteq G$ is finite, then 
\begin{align*}
  \mu\left( A \right) &= \frac{1}{\left\vert F \right\vert}\sum_{s\in F} \mu\left( sA \right)\\
                      &\leq \frac{1}{\left\vert F \right\vert}.
\end{align*}
To reframe this, we can consider an enumeration $s_0,s_1,s_2,s_3,\dots$ where $s_0 = e$, and consider the map $f\colon A\rightarrow X$ given by $f\left( s_n x \right) = s_{n+1}x$. This allows us to define a measure-preserving map $f$ from
\begin{align*}
  A' &= \bigcup_{n=0}^{\infty}s_nA
\end{align*}
to itself, with
\begin{align*}
  \mu\left( A \right) &= \mu\left( A'\setminus f\left( A' \right) \right)\\
                      &= \mu\left( A' \right)-\mu\left( f\left( A' \right)\right)\\
                      &= 0.
\end{align*}
This is a version of the pigeonhole principle that goes by the term Poincaré recurrence.
\begin{theorem}
  Let $G\curvearrowright \left( X,\mu \right)$ be a pmp action of an infinite group, and let $A$ be a subset with $\mu\left( A  \right)> 0$. Then, for almost every $x\in A$, the set of all $s\in G$ with $sx \in A$ is infinite.

  If $T\colon X\rightarrow X$ is a single pmp transformation, then for almost every $x\in A$, there are infinitely many $n\in \N$ such that $T^n(x) \in A$.
\end{theorem}
\begin{proof}
  Suppose not. Then, there is a finite $F\subseteq G$ such that
  \begin{align*}
    B &= A\setminus \bigcup_{s\in G\setminus F} sA
  \end{align*}
  has nonzero measure. Since $G$ is infinite, we may recursively construct an infinite set $I\subseteq G$ such that $s\notin tF$ For all distinct $s,t\in I$. The sets $sB$ for each $s\in I$ are then pairwise disjoint and have the same measure as $B$, which is a contradiction to the fact that $\left( X,\mu \right)$ is a finite measure space.

  In the case of a single transformation $T$, we find $n$ such that
  \begin{align*}
    A\setminus \bigcup_{k=n+1}^{\infty}T^{-k}A
  \end{align*}
  has nonzero measure and proceed in a similar manner.
\end{proof}
\section{Mixing and Weakly Mixing Actions}%
\begin{definition}
  The action $G\curvearrowright \left( X,\mu \right)$ is said to be mixing if $G$ is infinite and for all measurable $A,B\subseteq X$, the function
  \begin{align*}
    s\mapsto \mu\left( sA\cap B \right)-\mu\left( A \right)\mu\left( B \right)
  \end{align*}
  vanishes at infinity, in that for every $\ve > 0$, there is a finite subset $F\subseteq G$ such that the image of $G\setminus F$ is contained in the ball of radius $\ve$ centered at $0$.
\end{definition}
Since the independence of two measurable sets $A$ and $B$ is equivalent to the orthogonality of $\chi_A-\mu\left( A \right)1$ and $\chi_B - \mu\left( B \right)1$ in $L^2\left( X,\mu \right)\ominus \C 1$, we may rephrase this asymptotic independence condition to an asymptotic orthogonality property for unitary representations.
\begin{definition}
  A unitary representation $\pi\colon G\rightarrow B\left( H \right)$ is said to be \textit{mixing} if $G$ is infinite and, for all $\xi,\eta\in H$, the function $s\mapsto \iprod{\pi(s)\xi}{\eta}$ vanishes at infinity.
\end{definition}
These are equivalent when the particular unitary representation is the Koopman representation.
\begin{proposition}
  The action $G\curvearrowright \left( X,\mu \right)$ is mixing if and only if the restriction of the Koopman representation $\kappa$ to the orthogonal complement $L^2\left( X,\mu \right)\ominus \C 1$ of the constant functions is mixing.
\end{proposition}
\begin{proof}
  Suppose the action is mixing. Let $\xi,\eta\in L^2\left( X,\mu \right)\ominus \C 1$. We will show that $s\mapsto \iprod{\kappa(s)\xi}{\eta}$ vanishes at infinity. Via some bootstrapping, we may assume that there are finite partitions $\mathcal{P}$ and $\mathcal{Q}$ of $X$ and scalars $c_A$ and $d_B$ such that 
  \begin{align*}
    \xi &= \sum_{A\in \mathcal{P}} c_A\chi_A\\
    \eta &= \sum_{B\in \mathcal{Q}} d_B \chi_B.
  \end{align*}
  Then, we have
  \begin{align*}
    0 &= \int_{X}^{} \xi \overline{\eta}\:d\mu\\
      &= \sum_{A\in \mathcal{P}}\sum_{B\in \mathcal{Q}} c_A \overline{d_B} \mu\left( A \right)\mu\left( B \right),
  \end{align*}
  so by the triangle inequality, we have for all $s\in G$ that
  \begin{align*}
    \left\vert \iprod{\kappa(s)\xi}{\eta} \right\vert &\leq \sum_{A\in \mathcal{P}}\sum_{B\in \mathcal{Q}} \left\vert c_A \overline{d_B} \right\vert \left\vert \mu\left( sA\cap B \right)-\mu\left( A \right)\mu\left( B \right) \right\vert,
  \end{align*}
  which gives our asymptotic vanishing.

  In the reverse direction, we may apply the mixing hypothesis to the vectors $\chi_A-\mu(A)1$ and $\chi_B-\mu(B)1$ in $L^2\left( X,\mu \right)\ominus \C 1$.
\end{proof}
Mixing, while powerful, is a quite rare phenomenon. Therefore, we may weaken it in a suitable fashion as follows.

If $T\colon X\rightarrow X$ is a pmp transformation, then we define the weak mixing condition to be the mean asymptotic independence
\begin{align*}
  \lim_{n\rightarrow\infty} \frac{1}{n}\sum_{k=0}^{n-1} \left\vert \mu\left( T^{-k}A\cap B \right)-\mu\left( A \right)\mu\left( B \right) \right\vert &= 0.
\end{align*}
For finite subsets of a group $G$, the analogue of this condition is the Følner condition for amenability, and upon defining amenable actions in a suitable fashion, a similar definition for weak mixing holds.

Yet, there is a way to define this for a pmp action of any group, amenable or not.

Recall that if $f\in \ell_{\infty}\left( G \right)$, then there are two actions $\lambda_s(f)$ and $\rho_s(f)$ given by
\begin{align*}
  \lambda_s(f)(t) &= f\left( s^{-1}t \right)\\
  \rho_s(f)(t) &= f\left( ts^{-1} \right),
\end{align*}
which are called the left and right actions of $G$ respectively.
\begin{definition}
  A function $f\in \ell_{\infty}\left( G \right)$ is said to be weakly almost periodic if the set
  \begin{align*}
    \overline{\underbrace{\set{\lambda_s(f) | s\in G}}_{\eqcolon\lambda(G)(f)}}^{w}
  \end{align*}
  is weakly compact in $\ell_{\infty}\left( G \right)$. The set of weakly almost periodic functions is denoted $\operatorname{WAP}\left( G \right)$.
\end{definition}
\begin{proposition}
  The space $\operatorname{WAP}\left( G \right)$ is a unital $C^{\ast}$-subalgebra of $\ell_{\infty}\left( G \right)$ that is invariant under both $\lambda$ and $\rho$.
\end{proposition}
\begin{proof}
  By definition, $\operatorname{WAP}\left( G \right)$ is left invariant, contains the constant functions, and is closed under conjugation. Since the map $h\mapsto \rho_t(h)$ is bounded and linear, it is weakly continuous, so we have $ \overline{\rho_t \lambda(G)(f)}^{w} = \rho_t \left( \overline{\lambda(G)(f)}^{w} \right) $ by continuity, whence $\operatorname{WAP}\left( G \right)$ is right invariant.

  Furthermore, $\operatorname{WAP}\left( G \right)$ is linear, since if $\lambda\in \C$ and $f,g\in \operatorname{WAP}\left( G \right)$, then the map $\left( k,h \right)\mapsto \lambda k + h$ is weakly continuous on the product space of the weak orbit closures $ \overline{\lambda(G)(f)}^{w}\times \overline{\lambda(G)(g)}^{w} $, whence its image is weakly compact.

  Now, we show that $\operatorname{WAP}\left( G \right)$ is norm-closed. For this, the Eberlein--Smulian theorem gives that it suffices to show that for any $f$ in the norm closure of $\operatorname{WAP}\left( G \right)$, and any sequence $\left( s_k \right)_{k}$ in $G$, there is a weakly convergent subsequence $\left( \lambda_{s_k}(f) \right)_{k}\rightarrow f$. Using the Eberlein--Smulian theorem recursively and extracting the diagonal gives a sequence $\left( t_k \right)_k$ such that for each $n$, the sequence $\left( \lambda_{t_k}\left( f_n \right) \right)_{k}\rightarrow h_n\in \ell_{\infty}\left( G \right)$. Since the sequence $\left( h_n \right)_n$ is Cauchy in $\ell_{\infty}\left( G \right)$, there is some $h\in \ell_{\infty}\left( G \right)$. Given $\varphi$ in the unit ball of $\ell_{\infty}\left( G \right)^{\ast}$ and any $\ve > 0$, there is some $n$ such that $\norm{f-f_n} < \ve/3$ and $\norm{h_n - h} < \ve/3$. For sufficiently large $k$, we then have
  \begin{align*}
    \left\vert \varphi\left( \lambda_{t_k}\left( f_n \right) - h_n \right) \right\vert &< \ve/3
      \intertext{whence}
    \left\vert \varphi\left( \lambda_{t_k}\left( f \right) - h \right) \right\vert &< \ve,
  \end{align*}
  so that $\lambda_{t_k}\left( f_n \right)\rightarrow h$ weakly as $k\rightarrow\infty$, meaning $\operatorname{WAP}(G)$ is closed.

  Finally, we will show that $\operatorname{WAP}\left( G \right)$ is closed under multiplication. Let $f,g\in \operatorname{WAP}\left( G \right)$. Regarding $\ell_{\infty}\left( G \right) = C\left( \beta G \right)$, we have that the weak closure measures of $ \lambda(G)(f) $ and $\lambda(G)(g)$ are weakly sequentially compact. Therefore, by considering point masses on $\beta G$, given a sequence in $G$, we may extract a subsequence $\left( s_n \right)_n$ such that $\left( \lambda_{s_n}\left( f \right) \right)_{n}$ converges pointwise, and there is a subsequence $\left( \lambda_{s_{n_k}}\left( g \right) \right)_{k}$ that converges pointwise, so $\left( \lambda_{s_{n_k}}\left( fg \right) \right)_{k}\rightarrow h\in C\left( \beta G \right)$. By dominated convergence, it follows that
  \begin{align*}
    \int_{}^{} \lambda_{s_{n_k}}\left( fg \right)\:d\mu &\rightarrow \int_{}^{} h\:d\mu
  \end{align*}
  for all positive finite Radon measures $\mu$ on $\beta G$. Since the space of all positive Radon measures spans $C\left( \beta G \right)^{\ast}$, it follows that $\lambda_{s_{n_k}}\left( fg \right)\rightarrow h$ weakly, so $ \lambda(G)\left( fg \right) $ is weakly compact by Eberlein--Smulian, so $fg\in \operatorname{WAP}\left( G \right)$.
\end{proof}
The notion of weakly almost periodic functions is very useful in the study of weakly mixing representations.

In general, we know that it is not always the case that $\ell_{\infty}\left( G \right)$ admits an averaging function. Yet, there \textit{is} a unique $G$-invariant mean on $\operatorname{WAP}\left( G \right)$.
\begin{theorem}
  On $\operatorname{WAP}\left( G \right)$, there is a unique left-invariant mean, which coincides with a unique right-invariant mean.
\end{theorem}
\begin{proof}
  Let $f\in \operatorname{WAP}\left( G \right)$. By Krein--Smulian, the closed convex hull of a weakly compact subset of a Banach space is weakly compact, so since $ \overline{\lambda(G)(f)}^{w} $ is weakly compact, so is $ \overline{\conv}\left( \lambda(G)(f) \right) $. By the Ryll--Nardzewski fixed-point theorem, the left action of $G$ on $ \overline{\conv}\left( \lambda(G)(f) \right) $ has a fixed point $\lambda$, and similarly the right action of $G$ on $ \overline{\conv}\left( \rho(G)(f) \right) $ has a fixed point $\rho$, which are necessarily constant functions.

  Take nets $\left( L_{\alpha} \right)_{\alpha}$ of convex combinations of left translation operators $f\mapsto \lambda_s(f)$ and $\left( R_{\beta} \right)_{\beta}$ of convex combinations of right-translation operators $f\mapsto \rho_s(f)$ such that $L_{\alpha}(f)\rightarrow \lambda$ and $R_{\beta}(f)\rightarrow \rho$. Given $\ve > 0$, we have some $\alpha$ and $\beta$ such that $\norm{L_{\alpha}(f)-\lambda} < \ve/2$ and $\norm{R_{\beta}(f) - \rho} < \ve/2$ respectively. Since the left and right actions commute, we have
  \begin{align*}
    \norm{\lambda-\rho} &= \norm{R_{\beta}\lambda - L_{\alpha}\rho}\\
                        &\leq \norm{R_{\beta}\left( \lambda-L_{\alpha}(f) \right)} + \norm{L_{\alpha}\left( R_{\beta}(f)-\rho \right)}\\
                        &< \ve,
  \end{align*}
  so that $\lambda = \rho$ and $\lambda$ and $\rho$ are unique.

  We obtain a map $m\colon \operatorname{WAP}\left( G \right)\rightarrow \C$, where we identify the constants with $\C$. This map is homogeneous (with respect to scalar multiplication) since both $L_{\alpha}$ and $R_{\beta}$ are themselves homogeneous, is both left and right invariant, and has $m\left( 1 \right) = 1$ with $m\left( f \right)\geq 0$ whenever $f\geq 0$.

  Our goal now is to show that, given $f,g\in \operatorname{WAP}\left( G \right)$, we have $m\left( f + g \right) = m\left( f \right) + m\left( g \right)$. Let $\ve > 0$, and take a convex combination $L$ of left translation operators such that $\norm{L(f) - m(f)} < \ve/2$. Since $ \overline{\conv}\left( \lambda(G)\left( L(g) \right) \right) \subseteq \overline{\conv}\left( \lambda(G)(g) \right) $, and $m\left( L(g) \right)$ and $m\left( g \right)$ are unique in $ \overline{\conv}\left( \lambda(G)(L(g)) \right) $ and $ \overline{\conv}\left( \lambda(G)(g) \right) $ respectively, it follows that $m\left( L(g) \right) = m\left( g \right)$. We may thus find some convex combination $L'$ of left translation operators such that $\norm{L'L\left( g \right) - m\left( g \right)} < \ve/2$. This gives
  \begin{align*}
    \norm{L'L\left( f+g \right) - \left( m(f) + m(g) \right)} &\leq \norm{L' \left( L(f)-m(f) \right)} + \norm{L'L(g)-m(g)}\\
                                                              &< \ve,
  \end{align*}
  so since $m\left( f + g \right)$ is the unique fixed point for the left action in $ \overline{\conv}\left( \lambda(G)\left( f+g \right) \right) $, it follows that $m\left( f + g \right) = m\left( f \right) + m\left( g \right)$.
\end{proof}
\begin{definition}
  A pmp action $G\curvearrowright \left( X,\mu \right)$ is called \textit{weakly mixing} if
  \begin{align*}
    m\left( s\mapsto \left\vert \mu\left( sA\cap B \right) - \mu\left( A \right)\mu\left( B \right) \right\vert \right) &= 0
  \end{align*}
  for all measurable $A,B\subseteq X$. 
\end{definition}
\begin{proposition}
  Every weakly mixing pmp action is ergodic.
\end{proposition}
\begin{proof}
  The definition of weak mixing implies that for every $G$-invariant measurable $A\subseteq X$, we have $\mu\left( A \right)-\mu\left( A \right)^2 = 0$, so $\mu\left( A  \right) = 0$ or $\mu\left( A \right) = 1$.
\end{proof}
Suppose $G$ is a group, and $\pi\colon G\rightarrow B\left( H \right)$ is a unitary representation. For any $\xi,\zeta\in H$, the map $f_{\xi,\zeta}(s) = \iprod{\pi(s)\xi}{\zeta}$ is weakly almost periodic.

In the Koopman representation, orthogonality is the analogue of probabilistic independence in the case of the action of the group itself. Therefore, we have two options for expressing orthogonality in terms of the mean; that is, $m\left( f_{\xi,\zeta} \right) = 0$ or $m\left( \left\vert f_{\xi,\zeta} \right\vert \right) = 0$. The latter is what will allow us to define weak mixing.
\begin{definition}
  Let $\pi\colon G\rightarrow B\left( H \right)$ be a unitary representation. A vector $\xi\in H$ is called \textit{weakly mixing} for $\pi$ if $m\left( \left\vert f_{\xi,\xi} \right\vert \right) = 0$. The representation $\pi$ is called weakly mixing if every vector in $H$ is weakly mixing.
\end{definition}
\begin{proposition}
  Every weakly mixing unitary representation is ergodic.
\end{proposition}
\begin{proof}
  We observe that if $\pi\colon G\rightarrow B\left( H \right)$ is a unitary representation, and $\xi$ is a $G$-invariant vector in $H$, then $\left\vert f_{\xi,\xi} \right\vert$ takes on the value $\norm{\xi}^2$, so this is the mean for the representation. Since $m\left( \left\vert f_{\xi,\xi} \right\vert \right) = 0$, it follows that $\norm{\xi} = 0$, so $\xi = 0$.
\end{proof}
\begin{lemma}
  Let $\pi\colon G\rightarrow B\left( H \right)$ be a unitary representation. A vector $\xi\in H$ is called weakly mixing if and only if $m\left( \left\vert f_{\xi,\zeta} \right\vert \right) = 0$ for all $\zeta\in H$.
\end{lemma}
\begin{proof}
  Suppose $\xi$ is weakly mixing in $H$, and let $\zeta\in H$. To show that $s\mapsto \left\vert \iprod{\pi(s)\xi}{\zeta} \right\vert$ has mean zero, we may assume that $\zeta$ has zero component in $\set{\pi(s)\xi}_{s\in G}^{\perp}$. Approximating where necessary, we may assume that $\zeta = \pi(t)\xi$ for some $t\in G$. Since $M$ is $G$-invariant, we have
  \begin{align*}
    m\left( \left\vert f_{\xi,\zeta} \right\vert \right) &= m\left( t\left\vert f_{\xi,\xi} \right\vert \right)\\
                                                         &= m\left( \left\vert f_{\xi,\xi} \right\vert \right)\\
                                                         &= 0.
  \end{align*}
\end{proof}
We also get the following proposition.
\begin{proposition}
  A unitary representation $\pi\colon G\rightarrow B\left( H \right)$ is weakly mixing if and only if $m\left( \left\vert f_{\xi,\zeta} \right\vert \right) = 0$ for all $\xi,\zeta\in H$.
\end{proposition}
We will soon show that the representation is weakly mixing if and only if the restriction of the Koopman representation to $L^2\left( X,\mu \right)\ominus \C 1$ is weakly mixing.
\begin{theorem}
  Let $\pi\colon G\rightarrow B\left( H \right)$ be a unitary representation. Then, $m\left( f_{\xi,\zeta} \right) = \iprod{P\xi}{\zeta}$, where $P$ is the orthogonal projection of $H$ onto the closed subspace of $G$-invariant vectors.
\end{theorem}
\begin{proof}
  The closed subspace of $G$-invariant vectors is the orthogonal complement of the set of vectors of the form $\pi(t)\eta-\eta$ for some $\eta\in H$ and $t\in G$. This follows from the fact that if $\zeta$ is a $G$-invariant vector, then for every $\eta\in H$ and $s\in G$, we have
  \begin{align*}
    \iprod{\pi(s)\eta - \eta}{\zeta} &= \iprod{\eta}{\pi\left( s^{-1} \right)\zeta-\zeta}\\
                                     &= 0,
  \end{align*}
  and if $\zeta$ is orthogonal to every vector of the form $\pi(t)\eta-\eta$, then
  \begin{align*}
    \iprod{\pi(s)\zeta-\zeta}{\eta} &= \iprod{\zeta}{\pi\left( s^{-1} \right)\eta-\eta}\\
                                    &= 0,
  \end{align*}
  so that $\pi(s)\zeta = \zeta$. Therefore, it suffices to assume that $\zeta$ is either $G$-invariant or of the form $\pi(s)\eta-\eta$. In the first case, we have
  \begin{align*}
    f_{\xi,\zeta}\left( s \right) &= \iprod{\xi}{\pi\left( s^{-1} \right)\zeta}\\
                                  &= \iprod{\xi}{\zeta},
  \end{align*}
  so that $m\left( f_{\xi,\zeta} \right) = \iprod{\xi}{\zeta}$, while in the second case, we have
  \begin{align*}
    m\left( f_{\xi,\zeta} \right) &= m\left( \lambda_s\left( f_{\xi,\eta} \right)-f_{\xi,\eta} \right)\\
                                  &= 0.
  \end{align*}
\end{proof}
\begin{definition}
  Let $\pi\colon G\rightarrow B\left( H \right)$ be a unitary representation. We say a vector $\xi\in H$ is compact if $ \overline{\pi(G)\xi} $ is compact. The representation is said to be compact if every vector in $H$ is compact.
\end{definition}
\begin{definition}
  A subset $K\subseteq G$ is said to be syndetic if there exists a finite subset $F\subseteq G$ such that $FK = G$. The set is called thickly syndetic if, for every nonempty finite $F\subseteq G$, $ \bigcap_{s\in G}sK $ is syndetic.
\end{definition}
\begin{proposition}
  A positive function $f\in \operatorname{WAP}\left( G \right)$ has $m\left( f \right) = 0$ if and only if for every $\ve > 0$, the set $f^{-1}\left( \left[ 0,\ve \right) \right)$ is thickly syndetic.
\end{proposition}
\begin{theorem}
  Let $\pi\colon G\rightarrow B\left( H \right)$ be a unitary representation. The following are equivalent;
  \begin{enumerate}[(i)]
    \item $\pi$ is weakly mixing;
    \item for every finite $\Omega\subseteq H$ and every $\ve > 0$, the set of all $s\in G$ such that $ \left\vert \iprod{\pi(s)\xi}{\zeta} \right\vert < \ve $ for all $\xi,\zeta\in \Omega$ is thickly syndetic;
    \item for every finite $\Omega\subseteq H$ and $\ve > 0$, there is $s\in G$ such that $ \left\vert \iprod{\pi(s)\xi}{\zeta} \right\vert < \ve $ for all $\xi,\zeta\in \Omega$;
    \item the only compact vector in $H$ is the zero vector;
    \item $\pi$ has no nonzero finite-dimensional subrepresentations;
    \item $\pi\otimes\rho$ is ergodic for every unitary representation $\rho$ of $G$;
    \item $\pi\otimes\rho$ is weakly mixing for every unitary representation $\rho$ of $G$;
    \item $\pi\otimes \overline{\pi}$ is ergodic.
  \end{enumerate}
\end{theorem}
\begin{proof}
  For (i) implies (ii), we use the fact that $m\left( \left\vert f_{\xi,\zeta} \right\vert \right) = 0$ for all $\xi,\zeta\in H$, so it holds for any finite subset, and we use the proposition. The implication (ii) implies (iii) follows immediately.

For (iii) implies (iv), let $\xi$ be a nonzero vector in $H$. Applying (iii) recursively, we construct a sequence $\left( s_n \right)_n$ in $G$ by setting $s_1 = e$ and $s_n$ such that we have
\begin{align*}
  \left\vert \iprod{\pi\left( s_n \right)\xi}{\pi\left( s_k \right)\xi} \right\vert < \frac{1}{2}\norm{\xi}^2
\end{align*}
for all $k = 1,\dots,n-1$. Hence, we get
\begin{align*}
  \norm{\pi\left( s_n \right)\xi-\pi\left( s_k \right)\xi}^2 &= \iprod{\pi\left( s_n \right)\xi-\pi\left( s_k \right)\xi}{\pi\left( s_n \right)\xi-\pi\left( s_k \right)\xi}\\
                                                             &= 2\norm{\xi}^2 - 2\re\left( \iprod{\pi\left( s_n \right)\xi}{\pi\left( s_k \right)\xi} \right)\\
                                                             &\geq 2\left( \norm{\xi}^2 - \left\vert \iprod{\pi\left( s_n \right)\xi}{\pi\left( s_k \right)\xi} \right\vert \right)\\
                                                             &\geq \norm{\xi}^2.
\end{align*}
Hence, the set $\set{\pi\left( s_n \right)\xi}_{n\in \N}$ is not totally bounded, so $\xi$ is not compact for $\pi$.

For (iv) implies (v), recall that any closed and bounded subset of a finite-dimensional Hilbert space is compact.

For (v) implies (vi), we let $\rho\colon G\rightarrow B\left( K \right)$ be a unitary representation, and suppose there is a nonzero vector $\xi\in H\otimes H$ such that $\left( \pi\otimes\rho \right)(s)\xi = \xi$ for all $s\in G$. There is a Hilbert--Schmidt operator $T\colon \overline{K}\rightarrow H $ such that $\pi(s) T \overline{\rho}(s)^{\ast} = T$, with $T^{\ast} = \overline{\rho}(s)T^{\ast}\pi(s)^{\ast}$. Therefore, we have $TT^{\ast}\pi(s) = \pi(s)TT^{\ast}$. Since $TT^{\ast}$ is a nonzero compact operator, it has a nonzero eigenvalue with finite-dimensional eigenspace $E_{\lambda}$. For any $\zeta\in E_{\lambda}$ and $s\in G$, we have
\begin{align*}
  TT^{\ast}\pi(s)\zeta &= \pi(s)TT^{\ast}\zeta\\
                       &= \lambda\pi(s)\zeta,
\end{align*}
so that $\pi(s)\zeta\in E_{\lambda}$, contradicting the fact that there are no nonzero finite-dimensional subrepresentations. Therefore, $\pi\otimes\rho$ must be ergodic. Since $ \overline{\pi} $ is a unitary representation, this immediately gives (viii).

For (viii) implies (i), given $\xi\in H$, we have
\begin{align*}
  m\left( \left\vert f_{\xi,\xi} \right\vert^2 \right) &= m\left( s\mapsto f_{\xi,\xi}(s) \overline{f_{\xi,\xi}(s)} \right)\\
                                                       &= m\left( s\mapsto \iprod{\pi(s)\xi}{\xi} \iprod{ \overline{\pi}(s)\xi }{\xi} \right)\\
                                                       &= m\left( s\mapsto \iprod{\left( \pi\otimes \overline{\pi} \right)(s)\xi}{\xi} \right)\\
                                                       &= 0,
\end{align*}
so Cauchy--Schwarz gives
\begin{align*}
  m\left( \left\vert f_{\xi,\xi} \right\vert \right)^2 &\leq m\left( \left\vert f_{\xi,\xi} \right\vert^2 \right)m\left( 1 \right)\\
                                                       &= 0,
\end{align*}
giving (i).

For (vii) implies (vi), we use the fact that every weakly mixing representation is ergodic. For (vi) implies (vii), let $\rho$ be a unitary representation of $G$. For every unitary representation $\sigma$ of $G$, we have $\pi\otimes \left( \rho\otimes\sigma \right)$ is ergodic, so $ \left( \pi\otimes\rho \right)\otimes \sigma $ is ergodic, so $\pi\otimes\rho$ is weakly mixing since (vi) implies (i).
\end{proof}
\begin{theorem}
  Every unitary representation of $G$ decomposes uniquely into a direct sum of a weakly mixing subrepresentation and a compact subrepresentation. Moreover, a unitary representation is compact if and only if it decomposes into a direct sum of finite-dimensional subrepresentations, which we may take to be irreducible.
\end{theorem}
\begin{proof}
  Let $\pi\colon G\rightarrow B\left( H \right)$ be a unitary representation. The compact vectors form a closed $G$-invariant subspace $H_{k}\subseteq H$, and similarly, the weakly mixing vectors form a closed $G$-invariant subspace $H_{m}\subseteq H$. From the equivalent of (i) and (iv) in the previous theorem, it follows that $H_m\perp H_k$, and if the orthogonal complement of $H_{m}\oplus H_k$ is nonzero, then it must contain a nonzero compact vector, In particular, this means $H = H_m\oplus H_k$.

  Since closed balls in finite-dimensional Hilbert spaces are compact, it follows that a direct sum of finite-dimensional unitary representations is compact. Conversely, if $\pi\colon G\rightarrow B(H)$i s a compact unitary representation, then we let $\mathcal{C}$ be the collection of sets of pairwise orthogonal finite-dimensional subrepresentations of $\pi$. By Zorn's Lemma, there is a maximal element $\Omega$, and $\bigoplus_{\rho\in\Omega}\rho$ must equal $\pi$ as else, its orthogonal complement would contain a finite-dimensional representation by the equivalence of (iv) and (v) in the previous theorem. We may then decompose each finite-dimensional representation into irreducibles by recursively splitting along $G$-invariant subspaces, a process which terminates by finite-dimensionality.
\end{proof}
We can now translate the criteria for pmp actions of groups to the Koopman representation restricted to $L^2\left( X,\mu \right)\ominus \C 1$, which will require a certain level of care. For brevity, we will write $\mu(f) = \int_{X}^{} f\:d\mu$.
\begin{theorem}
  For a pmp action $G\curvearrowright \left( X,\mu \right)$, the following are equivalent:
  \begin{enumerate}[(i)]
    \item the action is weakly mixing;
    \item the restriction of the Koopman representation to $L^2\left( X \right)\ominus \C 1$ is weakly mixing;
    \item for all $f,g\in L^2\left( X \right)$, one has $m\left( s\mapsto \left\vert \mu\left( \left( \lambda_s(f) \right)g \right) - \mu(f)\mu(g) \right\vert \right) = 0$;
    \item for all finite collections $\Omega$ of measurable subsets of $X$ and all $\ve > 0$, the set of all $s\in G$ such that 
      \begin{align*}
        \left\vert \mu\left( sA\cap B \right) - \mu\left( A \right)\mu\left( B \right) \right\vert < \ve
      \end{align*}
     for all $A,B\in \Omega$ is thickly syndetic;
    \item for every finite collection $\Omega$ of measurable subsets of $X$ and every $\ve > 0$, there exists $s\in G$ such that $\left\vert \mu\left( sA\cap B \right) - \mu\left( A \right)\mu\left( B \right) \right\vert < \ve$ for all $A, B\in \Omega$;
    \item the only compact elements in $L^2\left( X \right)$ under the Koopman representation are the almost everywhere constant functions;
    \item the restriction of the Koopman representation to $L^2\left( X,\mu \right)\ominus \C 1$ has no nonzero finite-dimensional subrepresentations;
    \item for every ergodic pmp action $G\curvearrowright \left( Y,\nu \right)$, the product action $G\curvearrowright \left( X\times Y,\mu\times\nu \right)$ is ergodic;
    \item the product action $G\curvearrowright \left( X\times X,\mu\times \mu \right)$ is ergodic.
  \end{enumerate}
\end{theorem}
\begin{proof}
  The implication (i) implies (iv) follows from the proposition showing that $m(f) = 0$ if and only if $f^{-1}\left( \left[ 0,\ve \right) \right)$ is thickly syndetic. The implication (iv) implies (v) follows immediately.

  For the implication (v) implies (vi), let $f\in L^2\left( X \right)$ be compact. To show that $f$ is constant almost everywhere, it suffices to show that if $D\subseteq \C$ is closed, then $A = f^{-1}\left( D \right)$ has measure either $0$ or $1$. Recursively applying (v), there is a sequence $\left( s_n \right)_n$ in $G$ satisfying
  \begin{align*}
    \lim_{n\rightarrow\infty} \mu\left( s_nA \cap s_m\left( X\setminus A \right) \right) &= \mu\left( A \right)\left( 1-\mu\left( A \right) \right)
  \end{align*}
  for all $m\in \N$ by selecting $s_1 = e$ and choosing $s_n$ to be such that for each $k=1,\dots,n-1$, we have  
  \begin{align*}
    \left\vert \mu\left( s_nA\cap s_k\left( X\setminus A \right) \right) - \mu\left( A \right)\mu\left( X\setminus A \right) \right\vert < 1/n
  \end{align*}
  Since $f$ is compact, we may assume that the sequence $\left( \lambda_{s_n}\left( f \right) \right)_{n}$ is Cauchy in the $L^2$ norm. For any $k\in \N$, let $C_k\subseteq X$ be the measurable subset of all $x\in X$ such that $f(x)$ is not within distance $1/k$ of the subset $D\subseteq \C$. Since $X\setminus A$ is equal to the union of the increasing sequence of subsets $C_k$ over $k$, then given some $\ve > 0$, there is some $k$ such that $\mu\left( X\setminus A \right) \leq \mu\left( C_k \right) + \ve$.

  For all $n,m$, this gives
  \begin{align*}
    \mu\left( s_nA\cap s_m\left( X\setminus A \right) \right) &\leq \mu\left( s_nA\cap s_mC_k \right) + \ve.
  \end{align*}
  By the definition of $C_k$, we have
  \begin{align*}
    \norm{\lambda_{s_n}\left( f \right)-\lambda_{s_m}\left( f \right)}_{L^2}^2 &\geq \frac{1}{k^2}\mu\left( s_nA\cap s_mC_k \right),
  \end{align*}
  so that $\lim_{m,n\rightarrow\infty}\mu\left( s_nA\cap s_mC_k \right) = 0$. Since $\ve$ may be arbitrarily small, we have $\mu_{m,n\rightarrow \infty} \mu\left( s_nA\cap s_m\left( X\setminus A \right) \right) = 0$, so we conclude that $\mu\left( A \right) = 0$ or $1$ as desired.

  For (vi) implies (vii), we again use the fact that any closed and bounded subset of a finite-dimensional Hilbert space is compact.

  For (vii) implies (viii), we let $G\curvearrowright \left( Y,\nu \right)$ be an ergodic pmp action. Then, $L^2\left( X\times Y \right)\ominus \C 1_{X\times Y}$ may be expressed as the orthogonal sum of $\left( L^2\left( X \right)\ominus \C 1_X \right)\otimes L^2\left( Y \right)$ and $\C 1_X\otimes \left( L^2\left( Y \right)\ominus \C 1_Y \right)$.

  The action on the first subspace is ergodic since the tensor product of a weakly mixing representation with any unitary representation is ergodic, and the action on the second subspace is ergodic since the action $G\curvearrowright \left( Y,\nu \right)$ is ergodic if and only if the Koopman representation restricted to $L^2\left( Y,\nu \right)\ominus \C 1_Y$ is ergodic. Therefore, $G$ acts ergodically on the direct sum.

  For the implication (viii) implies (ix), we let $Y$ be the trivial action on a one-point set, showing that $G\curvearrowright \left( X,\mu \right) $ is ergodic. Applying (viii) yields (ix) by letting the action on $Y$ be $G\curvearrowright \left( X,\mu \right)$.

  For the implication (ix) implies (ii), we observe that the Koopman representation $\kappa^2$ of $G\curvearrowright \left( X\times X,\mu\times \mu \right)$ is equal to the tensor product of the Koopman representation $\kappa$ via the canonical isomorphism $L^2\left( X\times X \right) \cong L^2\left( X \right)\otimes L^2\left( X \right)$. Since the unitary operator from $L^2\left( X \right)$ to $ \overline{L^2\left( X \right)} $ given by $f\mapsto \overline{f}$ intertwines $\kappa$ with the adjoint representation $ \overline{\kappa} $, we may identify $\kappa^2$ as $\kappa\otimes \overline{\kappa}$ such that the restriction of $\kappa^2$ to $L^2\left( X\times X \right)\ominus \C 1_{X\times X}$ contains the restriction of $\kappa$ to $L^2\left( X \right)\ominus \C 1_{X}$ tensored with its conjugate. This tensor product is ergodic by (ix), and so the restriction of $L^2\left( X \right)\ominus \C 1_X$ is weakly mixing.

  To obtain (ii) implies (iii), we observe that if $f,g\in L^2\left( X \right)$, then $f- \mu(f)1$ and $ \overline{g} - \mu\left( \overline{g} \right)1 $ are in $L^2\left( X \right)\ominus \C 1$, so that
  \begin{align*}
    m\left( s\mapsto \left\vert \iprod{sf - \mu\left( f \right)1}{ \overline{g} - \mu\left( \overline{g} \right)1 } \right\vert \right) &= 0.
  \end{align*}
  Expanding the inner product and simplifying gives the desired result.

  Finally, for (iii) implies (i), we apply (ii) to the indicator functions $\chi_A$ and $\chi_B$ for measurable subsets $A,B\subseteq X$.
\end{proof}
\begin{definition}
  A pmp action is said to be \textit{compact} if its Koopman representation is compact.
\end{definition}
\begin{theorem}
  Let $G\curvearrowright \left( X,\mu \right)$ be a pmp action. Then, the set $N$ of functions in $L^{\infty}\left( X \right)$ that are compact as elements of $L^2\left( X \right)$ under the Koopman representation is a $G$-invariant von Neumann subalgebra of $L^{\infty}\left( X \right)\subseteq B\left( L^2\left( X \right) \right)$, and the $L^2$ closure of $N$ in $L^2\left( X \right)$ is the subspace of compact vectors.
\end{theorem}
\begin{proof}
  The definition of compact vector gives that $N$ is a $G$-invariant linear subspace of $L^{\infty}\left( X \right)$ that is closed under the restriction of the $L^2$-norm topology.

  Now, let $f,g\in N$. We may assume that $f$ and $g$ have $L^{\infty}$ norm of at most $1$. Since $f$ and $g$ are compact, there is a finite $\Omega$ contained in the union of the orbits of $f$ and $g$ such that this union can be approximated to $\ve/2$ accuracy via $\Omega$. Then, for every $s\in G$, we can find $f',g'\in \Omega$ such that $\norm{\lambda_s(f)-f'}_{L^2} < \ve/2$ and $\norm{\lambda_s(g) - g'} < \ve/2$, which yields
  \begin{align*}
    \norm{\lambda_s\left( fg \right) - f'g'}_{L^2} &\leq \norm{\left( \lambda_{s}(f)-f' \right)\lambda_s(g)}_{L^2} + \norm{f'\left( \lambda_s(g)-g' \right)}_{L^2}\\
                                                   &\leq \norm{\lambda_s(g)}_{L^{\infty}}\norm{\lambda_s(f)-f'}_{L^2} + \norm{f'}_{L^{\infty}}\norm{\lambda_s(g)-g'}_{L^2}\\
                                                   &< \ve.
  \end{align*}
  Therefore, the set $\set{hk | h,k\in \Omega}$ is a finite subset that approximates the orbit of $fg$ to $\ve$ accuracy, meaning that the orbit is totally bounded. Therefore, $fg\in N$. Since $N$ is closed under taking complex conjugates, it follows that $N$ is a $G$-invariant von Neumann subalgebra of $L^{\infty}\left( X \right)$.
\end{proof}
\section{Important Examples}%
We will now focus on two of the most important examples of ergodic group actions: namely, the Bernoulli shift and the odometer. Both of these play an important role in various fields such as entropy theory and descriptive set theory.
\subsection{Bernoulli Shifts}%
Suppose $Y$ is a Polish space (for our purposes, a compact metrizable space or a finite set). If $G$ is a countable discrete group, then the product space $Y^G$ is Polish since $Y^G$ is separable and, if $d$ is a compatible metric on $Y$ with $\left( s_n \right)_n$ an enumeration of the elements of $G$, the metric
\begin{align*}
  d'\left( x,y \right) &= \set{n=1}^{\infty}2^{-n}\min\set{d\left( x_{s_n},y_{s_n} \right),1}
\end{align*}
is a compatible complete metric on $Y^G$.

The product topology on $Y^{G}$ is generated by the cylinder sets --- i.e., the sets $\prod_{s\in G}A_s$, where $A_s$ is open and $A_s = Y$ for all but finitely many $s\in G$. The open sets in $Y^G$ generate a Borel $\sigma$-algebra. If $\nu$ is a Borel probability measure on $Y$, then there is a unique Borel probability measure $\nu^G$ on $Y^G$ that acts as the product measure on cylinder sets; that is, if $F$ is a finite subset of $G$, and $\set{A_s}_{s\in F}$ is a collection of Borel subsets of $Y$, with $\pi_F\colon Y^G\rightarrow Y^F$ the coordinate restriction map, then
\begin{align*}
  \nu^G\left( \pi_{F}^{-1}\left( \prod_{s\in F}A_s \right) \right) &= \prod_{s\in F} \nu\left( A_s \right).
\end{align*}
If $G = \Z$, then we can view this as encoding probabilistic independence, in a sense. This measure can be constructed by defining the algebra $\mathcal{A}$ of cylinder sets, on which this definition of $\nu^{G}$ defines a premeasure. Then, Caratheodory's extension theorem yields existence and uniqueness of $\nu^{G}$.

We may define the action $G\curvearrowright Y^G$ by taking $\left( s\cdot x \right)_t = x_{s^{-1}t}$ for all $s,t\in G$ and $x\in Y^{G}$. This action preserves the measure $\nu^{G}$. In fact, it is a mixing action. To show this, we start exhibiting a measure theory fact.
\begin{proposition}
  Let $\left( X,\mu \right)$ be a probability space with $\sigma$-algebra $\mathcal{C}$, and let $\mathcal{A}$ be an algebra of measurable subsets of $X$ that generates $\mathcal{C}$ as a $\sigma$-algebra (save for null sets). Then, for every measurable $B\subseteq X$ and $\ve > 0$, there is $A\in \mathcal{A}$ such that $\mu\left( A\triangle B \right) < \ve$.
\end{proposition}
\begin{proof}
  By the hypothesis, $\mathcal{C}$ agrees with the smallest $\sigma$-algebra containing $\mathcal{A}$ save for null sets. Therefore, it suffices to show that the set $\mathcal{B}$ consisting of all measurable subsets $B$ satisfying the conclusion is a $\sigma$-algebra.

  We observe that $\emptyset$ belongs to $\mathcal{A}$, so it belongs to $\mathcal{B}$, and $\mathcal{B}$ is closed under complements, as $\left( X\setminus A \right)\triangle \left( X\setminus B \right) = A\triangle B$.

  Let $\left( B_n \right)_n\subseteq \mathcal{B}$ be a sequence, and let $\ve > 0$. Take $n\in \N$ such that
  \begin{align*}
    \mu\left( \bigcup_{k=n+1}^{\infty}B_k \right) &< \ve/2,
  \end{align*}
  and for each $k\in \N$ find $A_k\in \mathcal{A}$ such that $\mu\left( A_k\triangle B_k \right) < \ve/\left( 2n \right)$. Then,
  \begin{align*}
    \mu\left( \left( \bigcup_{k=1}^{n} \right)\triangle \left( \bigcup_{k=1}^{\infty}B_k \right) \right) &\leq \sum_{k=1}^{n}\mu\left( A_k\triangle B_k \right) + \mu\left( \bigcup_{k=n+1}^{\infty}B_k \right)\\
                                                                                                         &< \ve.
  \end{align*}
  Therefore, the union of all the $B_k$ lies in $\mathcal{B}$.
\end{proof}
If we take Borel cylinder sets $A = \prod_{s\in G}A_s$ and $B = \prod_{s\in G}B_s$ in $Y^G$, where $A_s$ and $B_s$ are equal to $Y$ outside of finite subsets $E$ and $F$ respectively, then $\mu\left( sA\cap B \right) = \mu\left( A \right)\mu\left( B \right)$ for all $s\in G\setminus FE^{-1}$. In particular, if $A$ and $B$ are finite unions of cylinder sets, then $\mu\left( sA\cap B \right) = \mu\left( A \right)\mu\left( B \right)$ for all $s$ outside a finite subset of $G$.

Since the cylinder sets form an algebra which generates the Borel $\sigma$-algebra, it follows that, if  $G$ is infinite, then $\lim_{s\rightarrow\infty}\mu\left( sA\cap B \right) = \mu\left( A \right)\mu\left( B \right)$ for all Borel $A,B\subseteq Y^{G}$. In particular, this means the Bernoulli actions of $G$ are mixing, hence weakly mixing and ergodic.

We can describe the Koopman representation $\kappa$ of $G\curvearrowright \left( Y^G,\nu^G \right)$ as follows. Take the infinite tensor product $L^2\left( Y \right)^{\otimes G}$, which is the completion of the inner product of the direct limit of the tensor products $L^2\left( Y \right)^{\otimes F}$ over finite subsets $F\subseteq G$. We determine the embeddings by taking $L^2\left( Y \right)^{\otimes E}\hookrightarrow L^2\left( Y \right)^{\otimes F}$ by defining $\bigotimes_{s\in E}\xi_s\mapsto \bigotimes_{s\in F} \widetilde{\xi}_s$, where $ \widetilde{\xi}_s = \xi_s $ if $s\in E$ and $ \widetilde{\xi}_s = 1 $ if $s\in F\setminus E$. We may then identify $L^2\left( Y^{G} \right)\cong L^2\left( Y \right)^{\otimes G}$.

Fixing an orthonormal basis $\Omega$ for $L^2\left( Y \right)$ containing $1$, and $\mathcal{C}$ the collection of all maps $\xi\colon G\rightarrow \Omega$ such that $\xi_s = 1$ for all but finitely many $s\in G$, then for any such $\xi$ we may write $\bigotimes_{s\in G}\xi_s$ for the corresponding vector in $L^2\left( Y \right)^{\otimes G}$ that is the elementary tensor $\bigotimes_{s\in F}\xi_s$ for the corresponding vector in $L^2\left( Y \right)^{\otimes F}$ whenever $F$ is a finite subset of $G$ with $\xi_s = 1$ for all $s\in G\setminus F$.

The set
\begin{align*}
  Z &= \set{\bigotimes_{t\in G}\xi_t : \xi\in \mathcal{C}}
\end{align*}
is then an orthonormal basis for $L^2\left( Y \right)^{\otimes G}$ which is invariant under $\kappa$, yielding an action $G\curvearrowright Z$ given by $s\cdot \bigotimes_{t\in G}\xi_t = \bigotimes_{t\in G}\xi_{s^{-1}t}$.

We may write this action as $G\curvearrowright \bigsqcup_{\zeta\in R}G/G_{\zeta}$, where $R$ is a collection of representatives of orbits and $G_{\zeta}$ is the stabilizer subgroup. The componentwise action $G\curvearrowright G/G_{\zeta}$ is given by $s\cdot \left( tG_{\zeta} \right) = stG_{\zeta}$. This gives that $\kappa$ decomposes as a direct sum $\bigoplus_{\zeta\in R}\lambda_{G/G_{\zeta}}$, where $\lambda_{G/G_{\zeta}}$ is the left \textit{quasiregular} representation on $\ell_2\left( G/G_{\zeta} \right)$ given by $\left( \lambda_{G/G_{\zeta}}(s)f \right)\left( tG_{\zeta} \right) = f\left( s^{-1}tG_{\zeta} \right)$. Whenever $\zeta = \bigotimes_{s\in G}1$, then $G_{\zeta} = G$, giving that the trivial representation is contained in the Koopman representation. Meanwhile, for every other $\zeta\in Z$, the stabilizer subgroup $G_{\zeta}$ is necessarily finite, since $\zeta$ appears in $L^2\left( Y \right)^{\otimes F}$ for some finite set $F\subseteq G$, giving that $G_{\zeta}\subseteq FF^{-1}$.

If $G$ has no nontrivial finite subgroups (such as $\Z$) and a nontrivial Bernoulli action, then we have that $\kappa$ is equivalent to $1 \oplus \lambda_{G}^{\oplus \N}$, where $1$ denotes the trivial representation and $\lambda_G$ the left regular representation. In the general case, the representation $\kappa$ is equivalent to a subrepresentation of $1 \oplus \lambda_G^{\oplus I}$, where $I$ is some countable index set, since if $H$ is a finite subgroup, the left quasiregular representation $\lambda_{G/H}$ is contained in the left regular representation $\lambda_G$ via the embedding $f\mapsto \left\vert H \right\vert^{-1/2}\left( f\circ q \right)$, where $q$ is the quotient map.

This enables us to show in the general case that whenever $G$ is infinite and the Bernoulli action is nontrivial, $\kappa$ is equivalent to $1\oplus \lambda_G^{\oplus \N}$. First, we show that $\kappa$ has a subrepresentation equivalent to $1\oplus \lambda_G^{\oplus \N}$.

Recursively, define finite subsets $F_1\subseteq F_2\subseteq\cdots$ by setting $t_1 = e$, $F_1 = \set{t_1}$, and for every $n > 1$, choosing $t_n\in G\setminus F_{n-1}F_{n-1}^{-1}F_{n-1}$, and setting $F_n = F_{n-1}\cup \set{t_n}$. Then, if there is $n \geq 3$ and $s\in G\setminus \set{e}$, we must have $sF_n \neq F_n$, as else we would have $1\leq i,j < n$ such that $st_i = t_j$ and $1\leq k < n$ with $st_n = t_k$, which would yield $t_n = t_it_j^{-1}t_k\in F_{n-1}F_{n-1}^{-1}F_{n-1}$. Pick $\eta\in \Omega\setminus \set{1}$, and choosing the vector $\bigotimes_{s\in G}\xi_{\eta,s}\in L^2\left( Y \right)^{\otimes G}$ where $\xi_{\eta,s} = \eta$ if $s\in F_n$ and $\xi_{\eta,s} = 1$ otherwise.

Each of these vectors generates a copy of the left regular representation, and since each of the $F_n$ have different cardinality, it follows that $\kappa$ contains a subrepresentation equivalent to $1 \oplus \lambda_G^{\oplus \N}$.

To prove that $\kappa\leq 1\oplus \lambda_G^{\oplus \N}$ and $1\oplus \lambda_G^{\oplus \N}\leq \kappa$ implies equivalence, we use some comparison theory of projections.
\begin{proposition}
  Let $\pi_1\colon G\rightarrow B\left( H_1 \right)$ and $\pi_2\colon G\rightarrow B\left( H_2 \right)$ be unitary representations such that each is equivalent to a subrepresentation of the other. Then, $\pi_1$ and $\pi_2$ are equivalent.
\end{proposition}
\begin{proof}
  Form the direct sum $\pi = \pi_1\oplus \pi_2$ on $H_1\oplus H_2$. Let $P_i\colon H_1\oplus H_2\rightarrow H_i$ be the orthogonal projection onto each subrepresentation. Define the operator $V$ in the commutant $\pi\left( G \right)'$ by composing $P_1$ with the isometry $H_1\rightarrow 0\oplus H_2$ that implements an equivalence between $\pi_1$ and a subrepresentation of $\pi_2$, and similarly define $W$ to be the composition of $P_2$ with the isometry $H_2\rightarrow H_1\oplus 0$ implementing the equivalence between $\pi_2$ and a subrepresentation of $\pi_1$. 

  Then, $V^{\ast}V = P_1$ and $VV^{\ast}\leq P_2$, while $W^{\ast}W = P_2$ and $WW^{\ast} \leq P_1$. Therefore, by the Cantor--Schröder--Bernstein theorem for projections, it follows there is $Z\in \pi\left( G \right)'$ such that $Z^{\ast}Z = P_1$ and $ZZ^{\ast} = P_2$, which yields an equivalence between $\pi_1$ and $\pi_2$.
\end{proof}
\subsection{Automorphisms of Compact Groups}%
Let $X$ be a compact group, and let $\mu$ be the normalized Haar measure on $X$. Let $G\curvearrowright X$ be an action by continuous group automorphisms. This action is necessarily $\mu$-preserving since the normalized Haar measure is unique.

When $X$ is abelian, we often refer to this as an \textit{algebraic action}. For instance, we may consider the $\Z$-action onto the $n$-torus $\R^n/\Z^n$ given by an element of $\SL_n\left( \Z \right)$. There is a general theory of algebraic actions that uses the \href{https://terrytao.wordpress.com/2011/09/27/254a-notes-3-haar-measure-and-the-peter-weyl-theorem/}{Peter--Weyl theorem} and the induced action of $G$ onto the dual $ \widehat{X} $,\footnote{In the general theory, this is the set of equivalence classes of irreducible unitary representations.} but we will focus on the abelian case.

If $A$ is a locally compact abelian group, then the Pontryagin dual $\widehat{A}$ is the group of all continuous homomorphisms $A\rightarrow S^1$ (characters) with the topology of uniform convergence on compact sets. \textit{Pontryagin duality} asserts that the map $A\rightarrow \widehat{\widehat{A}}$ taking $a\mapsto \left( \varphi\xmapsto{\hat{a}} \varphi(a) \right)$ is an isomorphism of topological groups. Similar to the general case of general locally compact groups, $\widehat{A}$ is the set of all equivalence classes of irreducible unitary representations (which are all one-dimensional since $A$ is abelian).

If $A$ is compact, then $\widehat{A}$ can be viewed as an orthonormal basis for $L^2\left( A,\mu \right)$. Observe that each element of $\widehat{A}$ is a common eigenvector for $\lambda(A)\subseteq B\left( L^2\left( A \right) \right)$, where $\lambda\colon A\rightarrow B\left( L^2\left( A \right) \right)$ is the left regular representation. Furthermore, $\lambda$ decomposes into a sum of the irreducible subrepresentations, which are given by elements of $\widehat{A}$.

If $G\curvearrowright X$ is an algebraic action, there is an action $G\curvearrowright \widehat{X}$ on the dual given by the automorphisms $s\cdot \varphi(x) = \varphi\left( s^{-1}\cdot x \right)$, where $\varphi\in \widehat{X}$, $x\in X$, and $s\in G$. Inverting this relationship by Pontryagin duality yields a one-to-one correspondence between algebraic actions $G\curvearrowright X$ and actions of $G$ on the discrete group $\widehat{X}$ by automorphisms. Actions of $G$ on $\widehat{G}$ gives rise to a left $\Z[G]$-module structure on $\widehat{X}$, so we obtain a one-to-one correspondence between algebraic actions of $G$ and left $\Z[G]$-modules.

Given an algebraic action $G\curvearrowright X$, the dual action $G\curvearrowright \widehat{X}$ corresponds to the Koopman representation $\kappa\colon G\rightarrow L^2\left( X,\mu \right)$, where we view $\widehat{X}$ as the orthonormal basis vectors of $L^2\left( X,\mu \right)$. The trivial subrepresentation of $X$, corresponding to the identity of $\widehat{X}$, acts on the constant functions in $L^2\left( X,\mu \right)$. If the action $G\curvearrowright X$ is not mixing, then by an approximation argument, there must be nonzero $\varphi,\rho\in \widehat{X}$ such that the set of all $s\in G$ such that $s\cdot \varphi = \rho$ is infinite. Denoting this set $K$, we then have that $s\cdot\varphi = \varphi$ for all $s\in K^{-1}K$.

This gives the following.
\begin{proposition}
  The action $G\curvearrowright X$ is mixing if and only if the stabilizer subgroup
  \begin{align*}
    G_{\varphi} = \set{s\in G : s\cdot\varphi = \varphi}
  \end{align*}
  is finite for every nontrivial $\varphi\in \widehat{X}$.

  In the special case of torsion-free $G$, then $G\curvearrowright X$ is mixing if and only if all these stabilizer subgroups are all trivial.
\end{proposition}
Given a character $\varphi\in \widehat{X}$, if the $G$-orbit of $\varphi$ is finite, then the linear span of this orbit in $L^2\left( X \right)$ is a finite-dimensional invariant subspace for the Koopman representation of $G$. Meanwhile, if the $G$-orbit of $\varphi$ is infinite, then every nonzero $\xi$ in the closed linear span of the orbit in $L^2\left( X \right)$ is not compact; by the decomposition of the Koopman representation into weakly mixing and compact parts, we observe that the weakly mixing subrepresentation acts on the closed linear span of the infinite $G$-orbits in $\widehat{X}\subseteq L^2\left( X \right)$, and the compact subrepresentation acts on the closed linear span of the finite $G$-orbits in $\widehat{X}\subseteq L^2\left( X \right)$.
\begin{proposition}
  If $G\curvearrowright X$ is an algebraic action, the following are equivalent:
  \begin{enumerate}[(i)]
    \item the action is ergodic;
    \item the action is weakly mixing;
    \item the orbit of every nontrivial element under the dual action $G\curvearrowright \widehat{X}$ is infinite.
  \end{enumerate}
\end{proposition}
\begin{proof}
  Observe that (ii) implies (i) immediately since it holds for general pmp actions. Similarly, (iii) implies (ii) follows from the observations above and the equivalence between (ii) and that the only compact elements are the (a.e.) constants.

  Finally, to see that (i) implies (iii), we observe that if $\varphi$ is a nontrivial element of $\widehat{X}$ with finite orbit $\set{\varphi_1,\dots,\varphi_n}$, then $\sum_{j=1}^{n}\varphi_j$ is nonzero and $G$-invariant in $L^2\left( X \right)\ominus \C 1$, which contradicts ergodicity.
\end{proof}
Any nontrivial subgroup of $\Z$ is infinite and has finite index; therefore, from the proposition above, any algebraic action of $\Z$ is such that ergodicity, weak mixing, and mixing are equivalent.
\subsection{Gaussian Actions}%
Given a Hilbert space $H$, we can consider its real portion to be a real Hilbert space $H_{\R}$ given by restriction of scalars. Similarly, for a unitary operator $u$ on $H$, the operator $u_{\R}$ is the orthogonal operator on $H_{\R}$ that is formally identical to $u$, and given a unitary representation $\pi\colon G\rightarrow B\left( H \right)$ of a group $G$, the orthogonal representation $\pi_{\R}\colon G\rightarrow B\left( H_{\R} \right)$ defined by $s\mapsto \pi(s)_{\R}$ is known as the realification of $\pi$.

If $H$ is a real Hilbert space, there is a complexification $H_{\C}$ of $H$ defined by taking pairs $\left( \xi,\eta \right)\in H\times H$, written as formal sums $\xi + \mathfrak{i}\eta$ (notice the fraktur here), and equipping them with the operations
\begin{align*}
  \left( \xi_1 + \mathfrak{i}\eta_1 \right) + \left( \xi_2 + \mathfrak{i}\eta_2 \right) &= \left( \xi_1 + \xi_2 \right) + \mathfrak{i}\left( \eta_1 + \eta_2 \right)\\
  \left( a+bi \right)\left( \xi + \mathfrak{i}\eta \right) &= \left( a\xi - b\eta \right) + \mathfrak{i}\left( b\xi + a\eta \right)\\
  \iprod{\xi_1 + \mathfrak{i}\eta_1}{\xi_2 + \mathfrak{i}\eta_2} &= \iprod{\xi_1}{\xi_2} + \iprod{\eta_1}{\eta_2} + i \iprod{\eta_1}{\xi_2} - i \iprod{\xi_1}{\eta_2}.
\end{align*}
An orthogonal operator $u$ on $H$ can then yield a unitary operator on $H_{\C}$ by taking $\xi + \mathfrak{i}\eta \mapsto u\xi + \mathfrak{i}u\eta$, which we denote $u_{\C}$. Similarly, an orthogonal representation $\pi\colon G\rightarrow B\left( H \right)$ of a group yields a unitary representation $\pi_{\C}\colon G\rightarrow B\left( H_{\C} \right)$ given by $s\mapsto \pi(s)_{\C}$, and is called the complexification of $\pi$.

Given a complex Hilbert space $H$, the maps $V,W\colon H\rightarrow \left( H_{\R} \right)_{\C}$ given by
\begin{align*}
  V\xi &= \frac{1}{\sqrt{2}} \left( \xi + \mathfrak{i}\left( -i\xi \right) \right)\\
  W\xi &= \frac{1}{\sqrt{2}} \left( \xi + \mathfrak{i}\left( i\xi \right) \right)
\end{align*}
yield linear and conjugate-linear isometries respectively into $\left( H_{\R} \right)_{\C}$. Furthermore, $W$ can be viewed as a Hilbert space embedding of the conjugate Hilbert space $ \overline{H} $, and if $\xi + \mathfrak{i}\eta$ is a vector in $\left( H_{\R} \right)_{\C}$, then
\begin{align*}
  \xi + \mathfrak{i}\eta &= \frac{1}{\sqrt{2}} V\left( \xi + i\eta \right) + \frac{1}{\sqrt{2}} W\left( \xi - i\eta \right).
\end{align*}
Furthermore, if $U$ is a unitary operator on $H$, then $VU = \left( U_{\R} \right)_{\C}V$, while $W \overline{U} = \left( U_{\R} \right)_{\C}W$, where $ \overline{U} $ is the formal equivalent operator on $ \overline{H} $ to $U$.
\begin{proposition}
  The map $\varphi\colon H\oplus \overline{H}\rightarrow \left( H_{\R} \right)_{\C}$ defined by
  \begin{align*}
    \varphi\left( \eta_1,\eta_2 \right) &= V\eta_1 + W\eta_2
  \end{align*}
  is a unitary isomorphism. Furthermore, if $\pi\colon G\rightarrow B\left( H \right)$ is a unitary representation, then $\varphi$ intertwines $ \pi \oplus \overline{\pi} $ with $\left( \pi_{\R} \right)_{\C}$.
\end{proposition}
The usefulness of this construction will emerge shortly. Namely, we will introduce the \textit{symmetric Fock space}, which will be used in our discussion of Gaussian Hilbert spaces. We will assume that the Hilbert space $H$ from which we are constructing this space is real.\footnote{In the complex case, we replace all instances of ``orthogonal'' with ``unitary'' and the base $H^{\odot 0}$ is defined to be $\C$ instead of $\R$.}

Let $n\in \N$. Define an orthogonal representation $\sigma\mapsto U_{\sigma}$ of $S_n$ onto the $n$th tensor power $H^{\otimes n}$ by taking
\begin{align*}
  \xi_1\otimes\cdots\otimes \xi_n &\mapsto \xi_{\sigma^{-1}\left( 1 \right)}\otimes\cdots\otimes \xi_{\sigma^{-1}\left( n \right)}.
\end{align*}
The $n$th symmetric power of $H$ is then the closed subspace
\begin{align*}
  H^{\odot n} &\coloneq \set{\xi\in H^{\otimes n} : U_{\sigma}\xi = \xi\text{ for all }\sigma\in S_n}.
\end{align*}
\begin{remark}
There are lots of connections here to representation theory of the symmetric group (i.e., Schur--Weyl duality), but we will not talk about those here.
\end{remark}
The orthogonal projection of $H^{\otimes n}$ onto $H^{\odot n}$ is given by
\begin{align*}
  P &= \frac{1}{n!} \sum_{\sigma\in S_n} U_{\sigma}.
\end{align*}
We write $f_1\odot\cdots\odot f_n$ for the image of an elementary tensor $f_1\otimes\cdots\otimes f_n$ under this projection.

If $H^{\odot 0}$ denotes $\R$, then the \textit{symmetric Fock space} $S\left( H \right)$ is then defined to be the Hilbert space direct sum $\bigoplus_{n=0}^{\infty}H^{\odot n}$.

If $U$ is an orthogonal operator on $H$, then the assignment
\begin{align*}
  \xi_1\odot\cdots\odot \xi_n &\mapsto U\xi_1\odot\cdots\odot U\xi_n
\end{align*}
on elementary tensors defines an orthogonal operator on $H^{\odot n}$ denoted $U^{\odot n}$. Then, we may define the orthogonal operator $S(U)$ on $S(H)$ by taking
\begin{align*}
  S(U) &= \bigoplus_{n=0}^{\infty}U^{\odot n}.
\end{align*}
The map $U\mapsto U^{\odot}$ from the orthogonal group of $H$ to the orthogonal group of $H^{\odot n}$ is a homomorphism, so every orthogonal representation $\pi\colon G\rightarrow B\left( H \right)$ gives rise to an orthogonal representation $s\mapsto \pi(s)^{\odot n}$ of $G$ on $H^{\odot n}$, denoted $\pi^{\odot n}$. The orthogonal representation $S(\pi)$ on $S(H)$ is given by
\begin{align*}
  S(\pi) &= \bigoplus_{n=0}^{\infty} \pi^{\odot n}.
\end{align*}
\begin{proposition}
  Let $H$ be a real Hilbert space. There is a canonical unitary isomorphism $S\left( H_{\C} \right)\rightarrow S\left( H \right)_{\C}$ which intertwines $S\left( \pi_{\C} \right)$ and $S\left( \pi \right)_{\C}$ for every orthogonal representation $\pi$ of a group on $H$.
\end{proposition}
\begin{proof}
  Define a map $\left( H_{\C} \right)^{\otimes n}$ to $\left( H^{\otimes n} \right)_{\C}$ by
  \begin{align*}
    \left( x_{1,0} + \mathfrak{i}x_{1,1} \right)\otimes\cdots\otimes \left( x_{n,0} + \mathfrak{i}x_{n,1} \right) &\mapsto \sum_{\omega\in \Omega} i^{\left\vert \omega^{-1}(1) \right\vert} x_{1,\omega(1)}\otimes\cdots\otimes x_{n,\omega(n)}\\
                                                                                                                        &+ \mathfrak{i}\sum_{\omega\in \Omega'} i^{\left\vert \omega^{-1}(1) \right\vert} x_{1,\omega(1)}\otimes\cdots\otimes x_{n,\omega(n)}.
  \end{align*}
  Here, $\Omega$ denotes the set of all functions $\set{1,\dots,n}\rightarrow \set{0,1}$ where $1$ is taken an even number of times, while $\Omega'$ is the set of all such functions that take $1$ an odd number of times.

  This is a unitary isomorphism that maps $\left( H_{\C} \right)^{\odot n}$ to $\left( H^{\odot n} \right)_{\C}$. Piecing together these isomorphisms along with the identity operator $\C\rightarrow \C$ for $n = 0$ yields isomorphisms $S\left( H_{\C} \right)\rightarrow S\left( H \right)_{\C}$.
\end{proof}
Next, we discuss Gaussian Hilbert spaces. This will require a crash course in probability theory.
\begin{definition}\hfill
  \begin{itemize}
    \item A \textit{random variable} on a probability space $\left( X,\mu \right)$ is a measurable function $f\colon X\rightarrow \R$.
    \item The \textit{distribution} of the random variable is the pushforward measure $f_{\ast}\mu \eqcolon \mu_f$. The \textit{joint distribution} of the measurable functions $f_1,\dots,f_k\colon X\rightarrow \R$ is the measure $\left( f_1,\dots,f_k \right)_{\ast}\mu$ on $\R^k$.
    \item When $f$ is integrable, we define
      \begin{align*}
        \mathbb{E}\left( f \right) &\eqcolon \int_{X}^{} f(x)\:d\mu(x)\\
                                   &= \int_{\R}^{} x\:d\mu_f(x).
      \end{align*}
      When $f\in L^2\left( X \right)$, then the \textit{variance} is defined by
      \begin{align*}
        \mathbb{V}\left( f \right) &= \mathbb{E}\left( \left( f- \mathbb{E}\left( f \right) \right)^2 \right)\\
                                   &= \mathbb{E}\left( f^2 \right) - \mathbb{E}\left( f \right)^2.
      \end{align*}
      We say $f$ is centered if $ \mathbb{E}(f) = 0 $.
    \item A family $\set{f_i}_{i\in I}$ of random variables is said to be independent if, for all finite subsets $F\subseteq I$ and Borel subset $B_i\subseteq \F$ with $i\in F$, the sets $f_i^{-1}\left( B_i \right)$ satisfy
      \begin{align*}
        \mu\left( \bigcap_{i\in F}f_i^{-1}\left( B_i \right) \right) &= \prod_{i\in F}\mu\left( f_i^{-1}\left( B_i \right) \right).
      \end{align*}
      In other words,
      \begin{align*}
        \mu_{\left( f_i \right)_{i\in F}} &= \prod_{i\in F} \mu_{f_i}
      \end{align*}
      for all finite $F\subseteq I$.
  \end{itemize}
\end{definition}
\begin{definition}
  A random variable $f$ on $\left( X,\mu \right)$ is said to be \textit{Gaussian} if it is either constant or its distribution $\mu_f$ satisfies
  \begin{align*}
    d\mu_f &= \frac{1}{\sqrt{2\pi \sigma^2}} e^{-\left( x-m \right)^2/2\sigma^2}\:d\lambda(x)
  \end{align*}
  for some $m\in \R$ and $\sigma > 0$, where $\lambda$ denotes the Lebesgue measure on $\R$.
\end{definition}
\begin{definition}
  A \textit{Gaussian Hilbert space} is a closed subspace $H$ of $L^2_{\R}\left( X \right)$ for some probability space $\left( X,\mu \right)$ such that the elements of $H$ are centered Gaussian random variables.
\end{definition}
Recall that the \textit{characteristic function} of a Borel probability $\mu$ on $\R^{k}$ is its Fourier transform (sans normalization constant)
\begin{align*}
  \left( t_1,\dots,t_k \right) &\mapsto \int_{}^{} e^{i\sum_{j=1}^{k}t_jx_j}\:d\mu\left( x_1,\dots,x_k \right).
\end{align*}
The characteristic function of a random variable is the characteristic function of ist distribution, i.e. the map
\begin{align*}
  \left( t_1,\dots,t_k \right) &\mapsto \mathbb{E}\left( e^{i\sum_{j=1}^{k}t_jf_j} \right)\\
                               &= \int_{}^{} e^{i\sum_{j=1}^{k}t_jx_j}\:d\mu_{\left( f_1,\dots,f_k \right)}\left( x_1,\dots,x_k \right).
\end{align*}
It can be shown that the following facts hold.
\begin{enumerate}[(i)]
  \item A Borel probability measure on $\R^k$ is uniquely determined by its characteristic function.
  \item If $f$ is a Gaussian random variable with mean $m$ and variance $\sigma^2$, then the characteristic function is given by $t\mapsto e^{imt-\sigma^2t^2/2}$.
  \item If $\set{f_1,\dots,f_n}$ is a family of independent random variables, then
    \begin{align*}
      \mathbb{E}\left( e^{itf_1}\cdots e^{itf_n} \right) &= \mathbb{E}\left( e^{itf_1} \right)\cdots \mathbb{E}\left( e^{itf_n} \right).
    \end{align*}
\end{enumerate}
\begin{proposition}
  If $\set{f_1,\dots,f_n}$ is an independent family of centered Gaussian random variables, then the linear span of $\set{f_1,\dots,f_n}$ is a Gaussian Hilbert space.
\end{proposition}
\begin{proof}
  If $\sigma_i^2$ denotes the variance of $f_i$, then for all $c_1,\dots,c_n\in \R$, the value of the characteristic function of the random variable $f = \sum_{j}c_jf_j$ at a given $t\in \R$ is
  \begin{align*}
    \mathbb{E}\left( e^{itf} \right) &= \mathbb{E}\left( e^{itc_1f_1}\cdots e^{itc_nf_n} \right)\\
                                     &= e^{-c_1^2\sigma_1^2t^2/2}\cdots e^{-c_n^2\sigma_n^2t^2/2}\\
                                     &= e^{-\left( c_1^2\sigma_1^2+\cdots+c_n^2\sigma_n^2 \right)t^2/2}.
  \end{align*}
  Therefore $f$ is a centered Gaussian random variable with variance $c_1^2\sigma_1^2 + \cdots + c_n^2\sigma_n^2$.
\end{proof}
\section{Multiple Recurrence and Szemeredi's Theorem}%
We know that if $\pi\colon G\rightarrow B\left( H \right)$ is a unitary representation, then this representation decomposes uniquely into an orthogonal direct sum of a weakly mixing representation and a compact representation; these representations are $G$-invariant.

If $G\curvearrowright \left( X,\mu \right)$ is a pmp action, then we can apply the decomposition to the Koopman representation of $G$ on $L^2\left( X \right)$, but this doesn't yield a description of the space $X$ as we may like. Fundamentally, $L^2\left( X \right)$ does not have the algebraic structure of $L^{\infty}\left( X \right)$, which is necessary to encode the dynamics of the action of $G$.

We start by developing the idea of a Hilbert module $L^2\left( X|Y \right)$ associated to a measure-preserving map $X\rightarrow Y$. Let $\left( X,\mu \right)$ and $\left( Y,\nu \right)$ be standard probability spaces, and let $\varphi\colon X\rightarrow Y$ be a measurable map such that $\nu = \varphi_{\ast}\mu$. We may regard $L^{\infty}\left( Y \right)\subseteq L^{\infty}\left( X \right)$ as a von Neumann subalgebra by the composition map $f\mapsto f\circ\varphi$.

Let $ \mathbb{E}_Y $ denote the orthogonal projection (or conditional expectation) $L^2\left( X \right)\rightarrow L^2\left( Y \right)$. Observe that for all $f\in L^{\infty}\left( X \right)$, we have
\begin{align*}
  \int_{X}^{} \mathbb{E}_Y\left( f \right)\:d\mu &= \iprod{ \mathbb{E}_Y\left( f \right) }{1}\\
                                                 &= \iprod{f}{ \mathbb{E}_Y\left( 1 \right) }\\
                                                 &= \iprod{f}{1}\\
                                                 &= \int_{X}^{} f\:d\mu.
\end{align*}
Moreover, if $\M_n\left( L^{\infty}\left( X \right) \right)$ and $\M_n\left( L^{\infty}\left( Y \right) \right)$ denote the $n\times n$ matrix algebras with entries in $L^{\infty}\left( X \right)$ and $L^{\infty}\left( Y \right)$ respectively, then
\begin{enumerate}[(i)]
  \item $ \mathbb{E}_Y $ maps $L^{\infty}\left( X \right)$ contractively onto $ L^{\infty}\left( Y \right) $;
  \item $ \mathbb{E}_Y $ is completely positive.
\end{enumerate}
For (i), if $f\in L^{\infty}\left( X \right)$, observe that for all $g,h\in L^{\infty}\left( Y \right)$, we have
\begin{align*}
  \left\vert \iprod{ \mathbb{E}_y\left( f \right)g }{h} \right\vert &= \left\vert \iprod{ \mathbb{E}_Y\left( f \right) }{h \overline{g}} \right\vert\\
                                                                    &= \left\vert \iprod{f}{ \mathbb{E}_Y\left( h \overline{g} \right) } \right\vert\\
                                                                    &= \left\vert \iprod{f}{h \overline{g}} \right\vert\\
                                                                    &= \left\vert \iprod{fg}{h} \right\vert\\
                                                                    &\leq \norm{f}_{L^{\infty}}\norm{g}_{L^2}\norm{h}_{L^2}.
\end{align*}
Therefore, we have
\begin{align*}
  \norm{ \mathbb{E}_Y\left( f \right)\left( \nu\left( D \right)^{-1/2}\chi_{D} \right) }_{L^2} &\leq \norm{f}_{L^{\infty}}
\end{align*}
for every set $D\subseteq Y$ with $\nu(D) > 0$. Therefore, $ \mathbb{E}_Y\left( f \right) $ is an element of $L^{\infty}\left( Y \right)$ with norm at most $\norm{f}_{L^{\infty}}$.
\nocite{kerr_and_li}
\printbibliography
\end{document}
